% This is the ADASS_template.tex LaTeX file, 24 Oct 2025.
% It is based on the ASP general author template file, but modified to reflect the specific
% requirements of the ADASS proceedings.

% This is the aspauthor.tex LaTeX file
% Copyright 2021, Astronomical Society of the Pacific Conference Series
% Revision:  29 November 2021

% To compile, at the command line positioned at this folder, type:
% pdflatex ADASS_template
% pdflatex ADASS_template
% This will create a file called ADASS_template.pdf

\documentclass[11pt,twoside]{article}

% Do NOT use ANY packages other than asp2023. 
\usepackage{asp2023}

\aspSuppressVolSlug
\resetcounters

% References must all use BibTeX entries in a .bibfile.
% References must be cited in the text using \citet{} or \citep{}.
% Do not use \cite{}.
% See ManuscriptInstructions.pdf for more details
\bibliographystyle{asp2021}

% The ``markboth'' line sets up the running heads for the paper.
% 1 author: "Surname"
% 2 authors: "Surname1 and Surname2"
% 3 authors: "Surname1, Surname2, and Surname3"
% >3 authors: "Surname1 et al."
% Replace ``Short Title'' with the actual paper title, shortened if necessary.
% Use mixed case type for the shortened title
% Ensure shortened title does not cause an overfull hbox LaTeX error
% See ASPmanual2010.pdf 2.1.4  and ManuscriptInstructions.pdf for more details
\markboth{Jankowsky, Aït Benkhali, and Wagner}{AI Assisted Observatory Management}

\begin{document}

\title{AI Assisted Observatory Management}

% Note the position of the comma between the author name and the 
% affiliation number.
% Authors surnames should come after first names or initials, eg John Smith, or J. Smith.
% Author names should be separated by commas.
% The final author should be preceded by "and".
% Affiliations should not be repeated across multiple \affil commands. If several
% authors share an affiliation this should be in a single \affil which can then
% be referenced for several author names. If only one affiliation, no footnotes are needed.
% See ManuscriptInstructions.pdf and ASP's manual2010.pdf 3.1.4 for more details
\author{F.~Jankowsky$^1$, F.~Aït~Benkhali$^1$, and S.~Wagner$^{1,2}$}
\affil{$^1$German Centre for Astrophysics (DZA), Görlitz, Germany; \email{felix.jankowsky@dzastro.de}}
\affil{$^2$Heidelberg State Observatory (LSW), Heidelberg, Germany}

% This section is for ADS Processing.  There must be one line per author. paperauthor has 9 arguments.
\paperauthor{Felix~Jankowsky}{felix.jankowsky@dzastro.de}{}{Deutsches Zentrum für Astrophysik}{Astrophysics}{Görlitz}{}{02826}{Germany}
\paperauthor{Faiçal~Aït~Benkhali}{faical.ait_benkhali@dzastro.de}{}{Deutsches Zentrum für Astrophysik}{Astrophysics}{Görlitz}{}{02826}{Germany}
\paperauthor{Stefan~Wagner}{swagner@lsw.uni-heidelberg.de}{}{Landessternwarte Heidelberg}{}{Heidelberg}{}{69117}{Germany}

% There should be one \aindex line (commented out) for each author. These are used to
% build up the author index for the Proceedings. The surname must come first, followed by
% initials. Note the use of ~ before each initial to control spacing.
% The \author entries (see above) have surname last. These \aindex entries have
% surname first.
% The Aindex.py command willl create them for you after you have constructed the \author
% The first entry should be the first author, for bold-facing the author index page-reference

%\aindex{Jankowsky,~F.}
%\aindex{Aït Benkhali,~F.}
%\aindex{Wagner,~S.}


\begin{abstract}
  In the past years the field of observatory operations has seen many technical advances which have predominantly been driven by the need for robotisation. Employing the new opportunities provided by AI agents advances this development. The German Centre for Astrophysics (DZA) has launched several projects in this regard. In this contribution we present (1) MeerKAT+RAG, which enables technical queries in a chatbot-like manner which bases information solely on existing documents, and (2) MARS, a Multi-Agent Review System that serves as support of observatory proposal evaluation.
\end{abstract}

% These lines show examples of subject index entries. At this stage these have to commented
% out, and need to be on separate lines. Eventually, they will be automatically uncommented
% and used to generate entries in the Subject Index at the end of the Proceedings volume.
% Don't leave these in! - replace them with ones relevant to your paper.

% These lines show examples of ASCL index entries. At this stage these have to commented
% out, and need to be on separate lines. Eventually, they will be automatically uncommented
% and used to generate entries in the ASCL Index at the end of the Proceedings volume.
% The ascl.py command will scan your paper on possible code names.
% Don't leave these in! - replace them with ones relevant to your paper.

\section{MeerKAT+RAG}

MeerKAT+RAG is a domain-specific retrieval-augmented generation (RAG) framework for answering technical and operational questions about the MeerKAT-Extension radio telescopes.
It makes answer verifiability a central architectural concern, transforming documentation into structured, provenance-aware chunks and combining dense semantic retrieval with lexical search via Reciprocal Rank Fusion (RRF), followed by cross-encoder reranking.
A deterministic knowledge graph and bounded agentic retrieval support entity-centred and multi-step queries.
Before delivery, answers undergo two-tier verification using deterministic claim and numeric checks, complemented by an optional language-model-based entailment assessment.
The framework can withhold answers when evidence is insufficient.
Evaluation on a 123-question benchmark informs the design and provides an initial assessment of its capabilities and limitations.
MeerKAT+RAG prioritises traceable evidence and explicit verification over fluent but potentially unsupported responses.

\subsection{Motivation}

MeerKAT-Extension commissioning operations depend on precise technical specifications and procedures, including wind limits, antenna pointing ranges, receiver frequency bands, cooldown procedures, equipment identifiers, and safety interlocks etc.
This knowledge is distributed across technical manuals and datasheets, embedded in prose, tables, lists, and figures.
Three challenges limit reliable access to this information:
\begin{itemize}
\item Lexical Search: can miss questions phrased differently from the source text 
\item Semantic retrieval: can overlook exact identifiers or confuse similar specifications and numerical values.
\item Unconstrained LLMs: can generate plausible but unsupported answers, particularly when evidence is incomplete
\end{itemize}

In telescope operations, a plausible answer is not necessarily a reliable one. Answers must be grounded in authoritative documentation, and insufficient evidence must remain a valid reason not to answer.
The goal was therefore to develop an end-to-end RAG workflow that makes MeerKAT-Extension technical documentation queryable while preserving source provenance and providing explicit mechanisms to verify answer support. 

\subsection{System Architecture}

MeerKAT+RAG separates offline knowledge preparation from online question answering.
The core retrieval pipeline is extended by a knowledge graph, bounded agentic search, figure-captioning capabilities, and answer-support checks.
These components are configuration-dependent rather than mandatory in every serving path.
Current configuration:
Qwen3-Embedding-4B for embeddings, Qwen3-Reranker-4B for reranking, and MiniMax-M3 for generation, accessed through an OpenAI-compatible endpoint.
These are implementation choices, not architectural requirements.

\begin{enumerate}
\item Ingestion and provenance: PDF/HTML sources are converted into Markdown with extracted tables, figures,and metadata. Document identity, revision, page context, and source-file hashes preserve links to the original material.
\item Structure-aware chunking: A section-aware chunker creates passages from prose, tables, lists, procedures, and code or configuration blocks. Stable chunk identifiers and document metadata support traceability across retrieval, agent logs, and evaluation.
\item Hybrid Indexing: Chunks are embedded and indexed in Qdrant, which also provides the full-text search used by the active lexical retrieval path. Dense semantic search captures paraphrases, while lexical matching helps retain exact identifiers and specifications. BM25 is used as an evaluation baseline, not as the production lexical implementation.
\item Evidence fusion and reranking: Reciprocal Rank Fusion (RRF) combines dense and full-text rankings without requiring their scores to be calibrated to a common scale. A cross-encoder reranks a bounded set of leading candidates, while the remaining candidates retain their RRF order.
\item Context-grounded generation: Selected passages and their source references are supplied to an OpenAI-compatible language model. Document, section, revision, content type, and page metadata support source tracing. Retrieval provides evidence for generation but does not guarantee correct interpretation.
\item Answer-support assessment: Two complementary checks are available
  \begin{itemize}
  \item Deterministic checks: examine textual overlap, numeric and unit-aware grounding, and reference validity. These checks are repeatable and inexpensive but cannot detect every semantic or logical error 
  \item LLM entailment assessment: an optional language-model judge provides an additional, advisory assessment of whether the evidence supports the answer
  \end{itemize}
\end{enumerate}

\subsection{Knowledge Graph and Agentic Retrieval}

\noindent \textbf{KG and Agent Loop:} A lightweight RDF knowledge graph, stored in SQLite, complements document retrieval through deterministic entity extraction and cue-anchored relations (seven node types and eight relation types). An optional agent combines five validated tools: graph-search, vector-search, hybrid-search, entity-extract, and lookup-chunk for multi-step retrieval, bounded to two decision iterations. Graph results use the same evidence representation as retrieved document chunks.

\noindent \textbf{Vision and Figure Understanding:} Integrated vision-model capabilities extend the framework beyond text retrieval by extracting information from figures and generating descriptions for downstream use.

\noindent \textbf{Interfaces and Operational Design:} MeerKAT+RAG supports CLI/Python, an OpenAI-compatible HTTP API, MCP, and Open WebUI. Stable chunk identifiers, recorded benchmark artefacts, validated tool inputs, bounded execution, and typed failures support traceability and reproducibility. These controls improve operational robustness but do not guarantee answer correctness.

\section{MARS}

MARS (Multi-Agent Review System) is an advisory AI framework that generates structured, evidence-grounded draft assessments for reducing the substantial workload of human reviewers.
Its five-role pipeline combines three parallel specialist analyses -- scientific merit, technical feasibility, and methodology/statistic -- with adversarial critique and consensus synthesis.
Substantive findings are linked to typed evidence, including proposal excerpts, verified literature records, and deterministic calculator outputs.
Unsupported claims are explicitly flagged.
A facility-configurable evaluation rubric is applied through a deterministic scoring layer separated from generative model judgments, while uncertainty and unresolved disagreements remain visible.
MARS supports, but does not replace, expert judgment:
It computes structured assessments and rubric scores.
Human reviewers retain full authority to inspect, revise, and finalize assessments.

\subsection{Motivation and Challenges}

Reviewers of astronomical proposals need to assess scientific merit, literature context, technical feasibility, observation times, and statistical robustness – while maintaining consistency and transparency. AI assisted review is promising, but generating a convincing review is not the same as producing a reliable one:
\begin{itemize}
\item Potential benefits: AI can support structured assessments, identify weaknesses, assist with first order review
\item Evaluation bias: LLM-based judges exhibit position, verbosity, and self-preference biases.
\item Security risks: AI agents using external tools can be vulnerable to malicious instructions in untrusted content. 
\item Coordination failures: AI systems can share blind spots, converge prematurely, or produce inconsistent findings
\end{itemize}

Turning promising AI capabilities into a reliable, secure, and auditable scientific review process requires explicit verification, uncertainty management, and human oversight.

\subsection{How MARS Addresses These Challenges}

How can AI-generated scientific assessments be made evidence-grounded, auditable, secure, and useful without displacing expert judgment?
MARS addresses this engineering challenge through five design objectives:
\begin{enumerate}
\item Evidence before assertion: Link substantive findings to inspectable proposal passages, verified literature records, or deterministic calculations; explicitly flag unverified claims. 
\item Structured independent analysis: Combine three specialized reviewers with adversarial critique and synthesis that preserves unresolved disagreements. 
\item Deterministic verification: Separate numerical calculations and rule-based score aggregation from generative model judgments.
\item Auditability and security: Record computational provenance, support replay and verification, and protect the pipeline against untrusted inputs and unsafe tool execution.
\item Human authority: Produce advisory draft reviews for expert inspection, correction, and final decision-making not autonomous allocation decisions.
\end{enumerate}

MARS shifts the focus from generating reviews to engineering a review process whose evidence, uncertainties, and computational history can be inspected.

\subsection{Evidence-Grounded Claims}

Every substantive claim is linked to typed evidence with an explicit verification status.
Unverifiable claims are flagged rather than presented as established facts.

\noindent \textbf{Proposal excerpts:} Exact text spans linked to the normalized proposal.

\noindent \textbf{Literature records:} Bibliographic metadata and source quotations verified against the original abstract.

\noindent \textbf{Calculator outputs:} Reproducible numerical checks of visibility, sensitivity, and exposure time.

\noindent \textbf{Unverified inferences:} Explicitly flagged assumptions or claims lacking verified support.

\noindent \textbf{Conflicting evidence:} Contradictory findings and their sources preserved for human review.

\noindent \textbf{Verification boundary:} Verifying a quotation does not verify its scientific interpretation. Evidence grounding improves traceability but does not guarantee correctness.

\subsection{Hierarchical Agent Workflow}

\begin{center}
  \begin{tabular}[t]{| c | c | c |}
    \hline
    \textbf{Agent/Role} & \textbf{Main responsibility} & \textbf{Typical output} \\
    \hline \hline
    Scientific & Assess novelty, theoretical soundness & Evidence-linked strengths, \\
    Reviewer & and braoder scientific impact & weaknesses and criterion \\
    & & scores \\
    \hline
    Technical & Examine observational feasibility & Feasibility finding and \\
    Reviewer & target visibility, instrument constraints & calculator records \\
    & sensitivity and exposure time & \\
    \hline
    Methodology & Examine analysis readiness, statistical & Findings and calculator \\
    Reviewer & assumptions, systematic uncertainties & record \\
    & and robustness & \\
    \hline
    Adversarial & Cross-check all three reports for & Critical report, no \\
    Critic & contradictions, unsupported claims & independent numerical \\
    & and overconfidence & scores \\
    \hline
    Consensus & Combine findings, respond to critique & Structured consensus \\
    Builder & and retain unresolved disagreement & report and advisory draft \\
    \hline
  \end{tabular}
\end{center}

\subsection{Scoring, Auditability and Human Oversight}

MARS combines structured rubric scoring, explicit uncertainty reporting, traceable execution, and human control.

\noindent \textbf{Explicitly six-criterion rubric:} Scientific novelty, theoretical soundness, observational feasibility, sensitivity adequacy, data-analysis readiness, and broader scientific impact.

\noindent \textbf{Transparent scoring:} Criterion-level scores and justifications are derived from structured agent findings using explicit rubric rules.

\noindent \textbf{Visible uncertainty:} Reports expose evidence links, unverified claims, missing information, critic warnings, unresolved disagreements, and measured run-to-run variability.

\noindent \textbf{Human oversight:} Reviewers can inspect evidence, edit scores and narratives, record overrides, and optionally assess proposals before seeing AI-generated drafts.

\noindent \textbf{Auditable provenance and replay:} Hash-linked records capture pipeline events, model calls, tool executions, scores, and human overrides. Portable run bundles support inspection and replay; discrepancies are flagged.

\noindent \textbf{Security and failure visibility:} Prompt-injection checks, restricted tools, sandboxing, and quarantine mitigate risks. Unverifiable citations, tool failures, and ambiguous calculations trigger explicit warnings.

% \bibliography{P101}  % For BibTex

% if we have space left, we might add a conference photograph here. Leave commented for now.
% \bookpartphoto[width=1.0\textwidth]{foobar.eps}{FooBar Photo (Photo: Any Photographer)}

\end{document}
